\chapter{Discussion}
\label{ch:discussion}

\section{What the Experiments Establish, and What They Do Not}
\label{sec:disc-claims}

The experiments support a narrow and conditional claim. Under one specific
configuration, at one dataset size, preferences constructed by sentence-level
aggregation are more consistently recoverable by a Bradley--Terry reward model
than preferences constructed holistically from the same evaluator and the same
candidate summaries. The difference averages roughly three percentage points over
six repeated runs and is consistent in direction across five of them.

Four claims that might appear to follow do not. The results do not establish that
GCA labels are factually more accurate, since no independent annotation was
collected and AlignScore supplies both sets of labels
(Section~\ref{sec:disc-reliability}). They do not establish that a summarization
model trained on GCA preferences produces more faithful summaries, since the final
study contains no policy-optimization stage. They do not establish that GCA is
generally preferable, since Chapter~\ref{ch:results} reports configurations and
scales at which it is not. And they do not establish a strong effect in absolute
terms, since both reward models remain near chance
(Section~\ref{sec:disc-weak-signal}). What the study contributes is evidence about
\emph{when} feedback granularity changes the structure of a preference-learning
signal, and the remainder of this chapter examines the conditions under which it
does.

\section{Why Granularity Interacts With Evaluator Design}
\label{sec:disc-interpretation}

The advantage reported in Chapter~\ref{ch:results},
Section~\ref{sec:results-final-campaign} is not a property of GCA in the abstract.
It appears only under a specific combination of two design choices,
simple-mean aggregation ($\alpha=0.0$) and the AlignScore \texttt{nli} evaluation
mode, and Section~\ref{sec:results-alpha0-validation} already showed that the
aggregation-formula change alone is not sufficient: at $\alpha=0.0$ under the
\texttt{nli\_sp} mode, holistic still led. This section asks why the
judge mode specifically makes the difference.

AlignScore's implementation \citep{zha2023alignscore} suggests a plausible
mechanism, though the present experiments do not test it directly. The \texttt{\_sp} suffix in \texttt{nli\_sp} enables
a preprocessing and aggregation step that is internal to AlignScore itself: the
context is split into roughly 350-token chunks and the claim is split into
sentences, and the final score is computed as, in the library's own
\texttt{inference\_per\_example} function, the maximum entailment score across
context chunks for each claim sentence, averaged across claim sentences
(\texttt{.max(dim=0).values.mean()}). Plain \texttt{nli} mode skips this
entirely and scores the full context against the full claim in one pass. The
chunk-then-aggregate pattern is not unique to AlignScore: SummaC applies the same
idea, segmenting document and summary and reducing a matrix of sentence-level
entailment scores to a single figure \citep{laban2022summac}. What differs
between evaluators is where that decomposition happens, which is precisely the
distinction this section turns on.

This has a direct consequence for what the holistic and GCA conditions actually
compute under each mode. Under \texttt{nli\_sp}, a holistic call already
decomposes the candidate summary into sentences internally, scores each one
against the best-matching chunk of the article, and averages the results, which
is functionally close to what GCA does \emph{externally} at $\alpha=0.0$: score
each summary sentence against the article and take the mean. When GCA is applied
on top of a judge that is already doing something close to sentence-level
decomposition and mean-pooling internally, its external decomposition adds little
genuinely new information, which is consistent with holistic and GCA performing
similarly, and even favoring holistic, under \texttt{nli\_sp}
(Table~\ref{tab:mode-sweep}). Under plain \texttt{nli}, by contrast, a holistic
call scores the entire summary against the (length-truncated) article as one
atomic unit, with no internal sentence decomposition at all. Here, GCA's external
per-sentence scoring is the only source of sentence-level localization in the
pipeline, and it is exactly under this mode that GCA shows a measurable
advantage.

The stored $n=200$ preference data allow one component of this account to be
checked directly rather than inferred. If a holistic \texttt{nli\_sp} call
internally splits the claim into sentences and averages, then the holistic score
of a summary should closely track the mean of the separately computed
per-sentence scores of that same summary. Across all 400 summaries in that set
the two quantities correlate at $r = 0.983$, agree to within $0.001$ for 234 of
400 summaries, and have a median absolute difference of $0.000$
(Figure~\ref{fig:mean-vs-holistic}). Under \texttt{nli\_sp}, holistic scoring and
mean-aggregated sentence scoring are therefore close to the same computation,
which is consistent with the observation that GCA at $\alpha=0.0$ adds little
under that mode.

\begin{figure}[htbp]
\centering
\begin{tikzpicture}
\begin{axis}[
  thesisplot,
  width=0.62\textwidth, height=6.2cm,
  xlabel={mean of per-sentence scores},
  ylabel={holistic score},
  xmin=0, xmax=1, ymin=0, ymax=1,
  axis equal image, grid=major,
]
\addplot[only marks, mark=*, mark size=1.0pt,
         draw=orange!70!black, fill=orange!55, fill opacity=0.55]
        table[x=meansent, y=holistic] {figures/mean_vs_holistic.dat};
\addplot[domain=0:1, samples=2, dashed, gray!70, line width=0.7pt] {x};
\end{axis}
\end{tikzpicture}
\caption[Holistic score vs.\ mean of per-sentence scores]{Holistic AlignScore against the mean of the separately computed
per-sentence scores, for all 400 candidate summaries in the $n=200$ preference
set under evaluator mode \texttt{nli\_sp}. The dashed line is $y=x$. The two
quantities are near-identical ($r = 0.983$), which is what the split mode's
internal sentence-averaging predicts. The corresponding files for the final
\texttt{nli} configuration are not retained in the repository, so the same check
could not be run for the unsplit mode.}
\label{fig:mean-vs-holistic}
\end{figure}

This measurement supports the part of the account concerning \texttt{nli\_sp},
but the account as a whole remains an interpretation of an exploratory sweep, and
three limitations constrain it. The mode comparison rests on a
single run per mode, and Section~\ref{sec:disc-stochasticity} shows that
single-run differences of this size fall within observed run-to-run variation;
only the selected \texttt{nli} configuration received repeated-run validation. No
ablation isolates the components of the \texttt{\_sp} preprocessing, such as chunk
size, from the mode change as a whole. And nothing here verifies empirically that
\texttt{nli\_sp}'s internal aggregation and GCA's external aggregation produce
comparable scores rather than merely comparable rankings. The proposition that
evaluator-internal decomposition reduces the marginal benefit of external GCA is
therefore offered as the most plausible reading of the observed pattern, and as a
hypothesis worth testing directly, rather than as a finding of this thesis
(Chapter~\ref{ch:conclusion}).

\section{How Much Weight the Statistical Evidence Can Carry}
\label{sec:disc-stats-caveat}

The $n=1{,}000$ result was first assessed from a fold-level interval and a
Wilcoxon $p$-value of 0.0034, computed over 30 pooled fold-level differences from
six runs. It is necessary to state plainly how much that figure could support, and
then to describe what changed once the campaign was extended.

Those 30 differences are not 30 independent observations. Within a single run the
five cross-validation folds partition one dataset: any two folds share four fifths
of their training data, and all five evaluate models trained from the same pool of
preference pairs. Their errors are correlated, so a test assuming independence
treats the evidence as stronger than it is. The effective sample size lies between
the number of runs and the number of folds, and closer to the former. The
fold-level $p$-value is therefore anti-conservative and is reported as an
exploratory summary, not as a hypothesis test, at either campaign size.

The appropriate unit of independent replication is the run. Applying the two-sided
Wilcoxon signed-rank test to the six run-level differences from the original
campaign ($+0.060$, $+0.013$, $-0.010$, $+0.051$, $+0.032$, $+0.039$) gives
$p = 0.0625$, which does not meet the conventional $0.05$ threshold. With six
observations, five of them positive, $0.0625$ is in fact the smallest two-sided
value the test can return, so that outcome reflected the number of runs available
at least as much as the size of the effect: no result with five wins and one loss
out of six paired observations can reach significance, regardless of how large or
small the true difference is.

This is a genuine limitation of a six-run design, not a property of the
underlying effect, and Chapter~\ref{ch:results},
Section~\ref{sec:results-extended-campaign} reports what happened when it was
addressed directly: the campaign was extended to twenty independently seeded
runs under a corrected, fully deterministic training procedure. GCA led in all
twenty, with a mean gap of $+3.95$ percentage points, a bootstrap 95\% confidence
interval of $[+3.19, +4.70]$ points that does not cross zero, and a run-level
two-sided Wilcoxon $p < 0.001$. This does meet the conventional significance
threshold, and the effect size is close to the original six-run estimate,
indicating that the earlier campaign's direction was not an artifact of having
too few runs to trust, only too few to prove.

Reaching significance changes what can be claimed, but it does not remove the
need for the qualifications that surrounded the six-run result. Statistical
significance at the run level establishes that GCA preferences are
\emph{consistently more learnable} than holistic preferences by this
reward-model architecture, under this evaluator configuration, at this dataset
size, to a degree unlikely to arise from run-to-run noise alone. It does not by
itself establish a large effect (Section~\ref{sec:disc-weak-signal} addresses
the absolute accuracy level), an effect that generalizes to other scales
(Section~\ref{sec:disc-scale}), or a mechanism (Section~\ref{sec:disc-interpretation}
remains an interpretation, not confirmed by the extended campaign). What the
evidence now supports is a directionally consistent, statistically significant
advantage of approximately three to four percentage points under this
configuration --- a materially stronger claim than the six-run campaign could
support, but one that is still scoped to the conditions under which it was
measured.

\section{Both Reward Models Are Weak in Absolute Terms}
\label{sec:disc-weak-signal}

A comparison between two conditions can be sound while both conditions perform
poorly, and that is the situation here. At $n=1{,}000$, the holistic reward model
reaches 0.5295 and the GCA reward model 0.5603 pairwise accuracy, against a chance
level of 0.5. In absolute terms these models recover very little of the preference
signal: the better of the two is correct on roughly 56 of every 100 held-out pairs,
where guessing would be correct on 50.

This materially qualifies what the thesis shows. The finding is that GCA
preferences are \emph{more learnable} than holistic preferences under the selected
configuration, not that either produces a reward model one would want to deploy.
A reward model at 0.56 accuracy would provide a very noisy training signal to any
downstream policy, and the significant $+3.95$ percentage-point advantage
established in the extended campaign (Chapter~\ref{ch:results},
Section~\ref{sec:results-extended-campaign}) is an improvement to a weak signal
rather than the difference between a weak and a strong one.

Several factors plausibly contribute to the low absolute numbers, and this thesis
can separate them only partially. The task itself is genuinely hard: both
candidate summaries in a pair are generated by the same model from the same
article and differ only in sampling temperature, so many pairs are near-ties in
factual quality, and a judge's preference between them may encode little
learnable structure. The RoBERTa-base reward model \citep{liu2019roberta} is also
small, and articles are truncated to 2{,}000 characters before tokenization
(Chapter~\ref{ch:implementation}, Section~\ref{sec:impl-rm-training}), so the model
may simply not see enough of the source to verify many claims. The accuracies at
$n=5{,}000$ and $n=10{,}000$ are higher for both conditions (0.58--0.59), which is
consistent with more training data helping the reward model in general, even as
the gap between conditions closes.

A preliminary ablation tested the truncation explanation directly by varying the
retained article length under the RoBERTa-base backbone, three seeds per setting
(Table~\ref{tab:truncation-ablation}). The result runs against the truncation
hypothesis as originally stated: accuracy was \emph{higher} with less retained
context (500 and 1{,}000 characters) than at the thesis's own 2{,}000-character
setting, for both conditions. A plausible reading is that 2{,}000 characters is
not actually more informative in practice, because the concatenation of article
prefix and summary already exceeds the 512-token cap at that length
(Chapter~\ref{ch:implementation}, Section~\ref{sec:impl-rm-training}), so how much
of the 2{,}000-character budget survives tokenization varies unpredictably with
summary length across examples. A shorter, more consistently-applied budget may
therefore produce a more uniform input distribution even though it discards more
raw text. This ablation does not yet isolate that explanation from a simple
backbone effect, since it was run only under RoBERTa-base; the two prepared but
unexecuted arms of the same ablation, under \texttt{deberta-v3-base} at 512 and
1{,}024 tokens, would hold the backbone fixed while varying context length and
are recorded as future work (Chapter~\ref{ch:conclusion}). Taken as it stands,
this result reframes rather than resolves the truncation question: raw article
length is not the limiting factor it was hypothesized to be, at least not in the
straightforward way of "more characters, more evidence, higher accuracy," and the
low absolute ceiling remains only partially explained.

\begin{table}[htbp]
\centering
\caption[Truncation ablation, RoBERTa-base arms]{Truncation ablation, RoBERTa-base arms (three seeds each). All three
rows use the 512-token cap; \texttt{max\_article\_chars} is the character budget
before tokenization. The \texttt{deberta-v3-base} arms (512 and 1{,}024 tokens)
were prepared but not completed; see Chapter~\ref{ch:conclusion}.}
\label{tab:truncation-ablation}
\begin{tabular}{rrrrr}
\toprule
Max article chars & Holistic & GCA & Gap (GCA $-$ Holistic) \\
\midrule
500   & 0.5553 & 0.5850 & +2.97 \\
1{,}000 & 0.5663 & 0.5903 & +2.40 \\
2{,}000 & 0.5290 & 0.5480 & +1.90 \\
\bottomrule
\end{tabular}
\end{table}

\section{Interpreting the Larger-Scale Results}
\label{sec:disc-scale}

The $n=1{,}000$ advantage was not reproduced with comparable magnitude in the
single runs at $n=5{,}000$ and $n=10{,}000$ (Chapter~\ref{ch:results},
Section~\ref{sec:results-scale}). Three things must be kept apart here: the
observed numerical behaviour, run-to-run variation, and hypotheses about
dataset-size effects. The observation is that two single runs produced gaps of
$-0.42$ and $+0.35$ percentage points where the pooled $n=1{,}000$ estimate was
$+3.08$. Whether this reflects a dataset-size effect is a separate question, and
the design does not answer it.

The best-supported explanation is statistical rather than substantive. The
$n=1{,}000$ estimate itself only stabilized after six pooled runs; the individual
per-run gaps in Table~\ref{tab:seed-validation}
ranged from $-0.010$ to $+0.060$, a spread wide enough that a single seed could
easily have landed on either side of zero. The $n=5{,}000$ and $n=10{,}000$
settings each contain a single run. Given how much the $n=1{,}000$ estimate moved
between its first two executions alone (from $+0.060$ down to $+0.013$), a single
run at a different dataset size landing outside the $n=1{,}000$ interval is not,
by itself, strong evidence that the underlying effect changes with dataset size.
It is at least as consistent with the same run-to-run variation already documented
at $n=1{,}000$, which simply went unmeasured at the larger sizes because only one
run was performed at each.

A substantive explanation is also possible, but this thesis does not have
evidence to support or rule it out: larger candidate pools could shift the
distribution of article lengths, sentence counts, or segmentation quality in ways
that change how much signal GCA's per-sentence scoring adds relative to holistic
scoring. This is stated here as an open hypothesis, not a finding, since nothing
in Chapter~\ref{ch:results} isolates article length, sentence count, or
segmentation quality as a variable.

Because these two explanations make different predictions, they are also
distinguishable in principle: if the scale effect is genuinely statistical noise,
repeated runs at $n=5{,}000$ should produce an interval overlapping the
$n=1{,}000$ result; if it reflects a real, scale-dependent change, they should
not. Repeating the $n=5{,}000$ experiment across several runs is the most direct
way to resolve this and is listed as future work in
Chapter~\ref{ch:conclusion}.

One hypothesis would predict this pattern, and it is worth stating even though the
present data cannot test it. At $n=1{,}000$ the reward
model is fitting a small dataset, and how cleanly the preference relation is
structured matters a great deal: a labeling that is internally more consistent is
easier to recover from few examples. As the dataset grows, the model sees more
examples of the same underlying relation, and additional data can compensate for a
noisier labeling. Under that account, the benefit of a cleaner preference
construction should be largest where data is scarce and should shrink as data
accumulates, which is the pattern observed. Consistent with this, absolute
accuracy rises for \emph{both} conditions between $n=1{,}000$ and $n=10{,}000$
(from 0.5295 and 0.5603 to 0.5827 and 0.5862), while the difference between them
contracts. The evidence is compatible with this explanation but does not establish
it, since scale and dataset composition vary together across these runs.

Read this way, the larger-scale results are informative rather than merely
disappointing. The observed pattern is consistent with the hypothesis that GCA may
provide a larger learnability advantage when preference data are limited, which
would also be the regime in which automatic preference construction is most
attractive, since it is where collecting human labels would be most burdensome.
The present design cannot establish such a boundary, because the larger-scale
settings contain single runs on nested datasets. What the study does contribute
here is the observation itself, together with a specification of the experiment
that would test it.

\section{Implications for RLAIF Preference Construction}
\label{sec:disc-implications}

Three implications follow for anyone constructing AI-generated preference data
from a fixed factuality judge, independent of whether the specific GCA method
studied here is adopted.

First, the aggregation formula used to combine sentence-level scores into a
preference label is not a minor implementation detail. The original
weakest-link-penalty formula ($\alpha=0.5$) actively hurt GCA relative to a plain
mean (Chapter~\ref{ch:results}, Table~\ref{tab:formula-comparison}), which runs
counter to the intuition that motivated it: a formula designed to be more
sensitive to localized errors was, in practice, more sensitive to noise than to
signal. Aggregation-formula choices of this kind should be validated
empirically rather than assumed correct from first principles. This is a
qualification worth attaching to the broader case for fine-grained supervision
\citep{wu2023finegrained}: locating errors more precisely helps only if the rule
that converts those locations into a training signal is itself sound.

Second, judge configuration can matter more than the granularity decision that
is nominally the point of comparison. Section~\ref{sec:disc-interpretation}'s
reading of AlignScore's mode behavior suggests that part of what looked like a
question about holistic-versus-sentence-level supervision was really a question
about how much sentence-level decomposition the judge was already doing
internally. Anyone comparing supervision granularities on top of a judge with
its own internal aggregation behavior should check what that judge already does
before attributing an effect to the granularity of their own pipeline. Surveys of
automatic judging emphasise sensitivity to prompt and configuration choices
\citep{zheng2023judging, gu2024survey}; the present result adds that a judge's
internal decomposition can also absorb the very manipulation an experiment
intends to study.

Third, a result validated at one dataset scale should not be treated as general
without a scale check. This is a caution this thesis can make credibly only
because Section~\ref{sec:disc-scale}'s robustness check was actually run rather
than assumed unnecessary; had the study stopped at the $n=1{,}000$ result, it
would have reported a materially less qualified conclusion than the evidence
supports.

\section{Judge Reliability Without Human Auditing}
\label{sec:disc-reliability}

No independent human validation of the factuality labels was collected in this
study, and it is worth stating precisely what this does and does not leave open.
The repeated runs in Chapter~\ref{ch:results},
Section~\ref{sec:results-final-campaign} establish that the difference between the
two preference sets is repeatedly observed: it survives changes in fold assignment
and training stochasticity, so it is not confined to a single favourable run.

Reproducibility and correctness are separate properties, however. That a
preference relation is consistently recoverable says nothing about whether the
underlying AlignScore judgments agree with what a human reader would consider
faithful to the source. A judge can produce labels that are highly reproducible
and systematically wrong, and nothing in this study would detect that case. The
labels function here as the reference standard, not as verified ground truth, and
every result is conditional on them in that sense. Establishing agreement with
human factuality judgments requires annotation this study did not collect,
quantified against a chance-corrected agreement measure such as Cohen's $\kappa$
\citep{cohen1960coefficient}, and the question is recorded as an open limitation
in Chapter~\ref{ch:conclusion}. The concern is not peculiar to this study:
reviews of reinforcement learning from human feedback identify the quality of the
preference data itself, and the difficulty of auditing it, as a central open
problem \citep{casper2023open}, and substituting an automatic judge for human
annotators relocates that problem rather than removing it.

\section{Stochastic Variation Between Runs}
\label{sec:disc-stochasticity}

Run-to-run variation is large enough in this setting to determine what conclusions
can be drawn, and it deserves treatment as a finding rather than as an
inconvenience. Two same-configuration runs differed by 4.7 percentage points
($+0.060$ versus $+0.013$), and across the six runs the gap ranged from $-0.010$
to $+0.060$. A study reporting only one of these runs could honestly have
concluded either that GCA helps substantially or that it does not help at all.

This is a known hazard rather than a peculiarity of the present setup. Fine-tuning
pretrained language models is sensitive to random seed alone, with weight
initialization and data order contributing comparably to the variance of
out-of-sample performance \citep{dodge2020finetuning}, and reporting score
distributions across several runs rather than a single figure has been urged for
precisely this reason \citep{reimers2017reporting}. The repeated-run design used
here follows that recommendation, and the spread it exposes is the reason this
thesis reports a distribution of gaps rather than a single headline number.

This variation traced back to an implementation property documented in
Chapter~\ref{ch:implementation}, Section~\ref{sec:impl-rm-training}: in the
original six-run campaign, the seed argument fixed the cross-validation fold
assignment through Python's standard random number generator, but PyTorch's
generator was not seeded in the cross-validation path. Several sources of
training stochasticity were therefore uncontrolled, including reward-head weight
initialization, the order in which training batches are drawn, and dropout
masks. Runs were consequently not exactly reproducible from a seed value alone.
This had an interpretive consequence at the time: the two executions at seed 42
were genuine repetitions rather than duplicates, because they shared the fold
assignment but not the training stochasticity, and the 4.7-point spread between
them reflects that uncontrolled stochasticity rather than fold assignment alone.

This property has since been fixed --- PyTorch and CUDA are now seeded per fold,
with an explicit generator for the training \texttt{DataLoader} and seeded
worker initialization, so a run is fully determined by its seed on fixed
hardware. The extended twenty-run campaign in Chapter~\ref{ch:results},
Section~\ref{sec:results-extended-campaign} was carried out under this corrected
procedure; each of its twenty seeds is consequently a genuinely independent and
exactly reproducible sample, which is what made it meaningful to extend the
campaign at all. The fix does not retroactively change the original six runs,
which remain a documented instance of exactly the seed-sensitivity the
literature describes, but it does mean the concern raised in this section no
longer applies to the results this thesis reports as its headline finding.

The practical implication generalizes beyond this study regardless of the fix.
At the accuracy levels and dataset sizes involved here, differences of one to two
percentage points between conditions are within the range that a single
under-seeded run can produce by chance. Comparisons of preference-construction
methods reported from single training runs at this scale, without controlling
every source of training randomness, should be treated with corresponding
caution.

\section{Confounds That Remain}
\label{sec:disc-confounds}

Two confounds could not be removed within this design, and both plausibly affect
the reported numbers.

The candidate pairs differ in decoding temperature, and
Chapter~\ref{ch:methodology}, Table~\ref{tab:temperature-asymmetry} shows the
consequence is not small: the lower-temperature candidate is preferred in 68.5\%
of holistic pairs and 63.0\% of GCA pairs on the data where this could be
measured. Because temperature systematically affects surface properties such as
length and lexical conservativeness, a reward model could recover a meaningful
share of the preference relation from stylistic cues without representing
factuality at all. The asymmetry also differs between conditions by roughly five
percentage points, so the two preference sets are not equivalent with respect to
this confound, and part of the observed difference between them could in principle
reflect that rather than granularity. Sampling both candidates at the same
temperature would remove the confound and is the cleanest correction available.

The reward model also sees only a truncated prefix of each source article
(Chapter~\ref{ch:implementation}, Section~\ref{sec:impl-rm-training}), while
subset selection admits articles several times longer than the retained prefix.
Where the evidence bearing on a claim lies beyond that prefix, the corresponding
preference is not recoverable from the model's input even in principle. This is a
plausible contributor to the low absolute accuracies discussed in
Section~\ref{sec:disc-weak-signal}, though its magnitude was not measured, and it
should not be assumed to be the dominant cause without the ablation described
there.

\section{Generalizability}
\label{sec:disc-generalizability}

Every result reported here is obtained under one combination of choices:
CNN/DailyMail as the corpus, AlignScore as the evaluator,
Mistral-7B-Instruct-v0.3 as the candidate generator, and a RoBERTa-base
Bradley--Terry model as the reward model. Each of these bounds the conclusions in
a distinct way.

The evaluator is the most consequential. Section~\ref{sec:disc-interpretation}
raises the possibility that whether external decomposition helps is related to
whether the evaluator already decomposes internally, which would make the finding
partly \emph{about} the evaluator's architecture rather than simply conditional on
AlignScore as a fixed instrument. If that reading is correct, a different
factuality model with different internal aggregation would shift the result. The corpus matters because CNN/DailyMail
summaries are short, extractive-leaning, and drawn from a single register; domains
with longer or more abstractive summaries would change both the number of
sentences per summary and the difficulty of the underlying factuality judgments.
The reward-model backbone matters because a stronger encoder, or one with a longer
context window, might recover more of the preference relation and thereby narrow
or widen the difference between conditions.

Taken together, these dependencies mean the study is best read as characterizing
behaviour under one specific configuration, and as motivating a hypothesis about
the interaction between external and evaluator-internal decomposition, rather than
as establishing a quantity that transfers directly to other setups.
